\begin{theorem}[abscissa powR le]
\uses{powR}
\label{abscissa_powR_le}
\lean{ArithmeticFunction.abscissa_powR_le}
% at /home/runner/work/trough/trough/library/PrimeNumberTheoremAnd/IwaniecKowalskiCh1.lean:595.0-608.55
\leanok
The abscissa of absolute convergence of $L(\text{pow}^R(\nu), s)$ is at most $\Re(\nu) + 1$.
\end{theorem}
\begin{proof}
\leanok
We apply \ref{LSeries.abscissaOfAbsConv_le_of_le_const_mul_rpow} which states that if there exists a constant $C$ such that $\|f(n)\| \leq C \cdot n^r$ for all $n$ sufficiently large, then the abscissa of absolute convergence of $L(f, s)$ is at most $r + 1$. In our case, we can take $f(n) = n^\nu$ and observe that $\|n^\nu\| = n^{\Re(\nu)}$. Thus, we can choose $C = 1$ and $r = \Re(\nu)$, which gives us the desired result that the abscissa of absolute convergence of $L(\text{pow}^R(\nu), s)$ is at most $\Re(\nu) + 1$.
\end{proof}
